%% thesis.tex 2014/04/11
%
% Based on sample files of unknown authorship.
%
% The Current Maintainer of this work is Paul Vojta.

\documentclass[masters]{ucbthesis}
\usepackage{amsmath}
\usepackage{latexsym}
\usepackage{amsfonts}
\usepackage[normalem]{ulem}
\usepackage{soul}
\usepackage{array}
\usepackage{amssymb}
\usepackage{extarrows}
\usepackage{graphicx}
\usepackage{hyperref}
\usepackage{framed}
\usepackage[utf8]{inputenc}
\usepackage[backend=biber,
natbib=true,
style=authoryear,
sorting=none,
isbn=true,
doi=true,
url=true,
]{biblatex}

\makeatletter
% special thanks to this guy
\newrobustcmd*{\parentexttrack}[1]{%
  \begingroup
  \blx@blxinit
  \blx@setsfcodes
  \blx@bibopenparen#1\blx@bibcloseparen
  \endgroup}

\AtEveryCite{%
  \let\parentext=\parentexttrack%
  \let\bibopenparen=\bibopenbracket%
  \let\bibcloseparen=\bibclosebracket}

\makeatother
\addbibresource{library.bib}

\usepackage{booktabs}
\usepackage{siunitx}
\usepackage{subfig}
\usepackage{wrapfig}
\usepackage{wasysym}
\usepackage{enumitem}
\usepackage{adjustbox}
\usepackage{ragged2e}
\usepackage[svgnames,table]{xcolor}
\usepackage{tikz}
% csv files and plots
\usepackage{pgfplots}
\pgfplotsset{compat=1.18}
\usepackage{csvsimple-l3}
\usepackage{longtable}
\usepackage{changepage}
\usepackage{setspace}
\usepackage{hhline}
\usepackage{multicol}
\usepackage{tabto}
\usepackage{float}
\usepackage{multirow}
\usepackage{makecell}
\usepackage{fancyhdr}
\usetikzlibrary{shapes.symbols,shapes.geometric,shadows,arrows.meta,
positioning,fit,backgrounds}
\tikzset{>={Latex[width=1.5mm,length=2mm]}}
\usepackage{tikz}
\usepackage[T1]{fontenc}
\usepackage{listings}

\usepackage{listings}
\lstset{
    language=C++, % Set the programming language
    backgroundcolor=\color{black!5}, % Light grey background
    basicstyle=\footnotesize\ttfamily, % Font style and size
    keywordstyle=\color{blue}, % Keyword color
    commentstyle=\color{green!60!black}, % Comment color
    stringstyle=\color{red}, % String color
    showstringspaces=false,
    numbers=left, % Line numbers on the left
    numberstyle=\tiny\color{gray}, % Line number style
    frame=single, % Add a frame around the code
    breaklines=true, % Automatic line breaking
    captionpos=b, % Caption position at the bottom
}


\urlstyle{same}

\renewcommand{\_}{\kern-1.5pt\textunderscore\kern-1.5pt}

%https://tex.stackexchange.com/questions/116534/lstlisting-line-wrapping
% I copied code from here to format things properly
% Credit to: Henri Menke who answered on StackExchange
\lstset{
  basicstyle=\ttfamily,
  columns=fullflexible,
  frame=single,
  breaklines=true,
  postbreak=\mbox{\textcolor{red}{$\hookrightarrow$}\space},
}

%%  Import package so that we can split the tex files into subfolders
% for easier compile and real-time edit

\usepackage{import}

%%%%%%%%%%%%%% for transfer function controllers %%%%%%%%%%%%%%
\usepackage{tikz}
\usetikzlibrary{shapes,arrows}
%%%% set citation style to square brackets


\begin{document}

% Declarations for Front Matter

\title{Derivations for Tampines Steam Tables and TAMPINES}
\author{Tom Lifesaver}
%\degreesemester{Spring}
%\degreeyear{1995}
%\degree{Doctor of Philosophy}
%\chair{Professor Richard Francis Sony}
%\othermembers{Professor Roger Spam \\
%  Associate Professor Michael Chex}
%% For a co-chair who is subordinate to the \chair listed above
%% \cochair{Professor Benedict Francis Pope}
%% For two co-chairs of equal standing (do not use \chair with this one)
%% \cochairs{Professor Richard Francis Sony}{Professor Benedict Francis Pope}
%\numberofmembers{3}
%% Previous degrees are no longer to be listed on the title page.
%% \prevdegrees{B.A. (University of Northern South Dakota at Hoople) 1978 \\
%%   M.S. (Ed's School of Quantum Mechanics and Muffler Repair) 1989}
%\field{Mathematics}
%% Designated Emphasis -- this is optional, and rare
%% \emphasis{Colloidal Telemetry}
%% This is optional, and rare
%% \jointinstitution{University of Western Maryland}
%% This is optional (default is Berkeley)
%% \campus{Berkeley}
%
%% For a masters thesis, replace the above \documentclass line with
%% \documentclass[masters]{ucbthesis}
%% This affects the title and approval pages, which by default calls this
%% document a "dissertation", not a "thesis".
%
%\maketitle
%% Delete (or comment out) the \approvalpage line for the final version.
%\approvalpage
%\copyrightpage

\include{abstract}

\begin{frontmatter}


% You can delete the \clearpage lines if you don't want these to start on
% separate pages.

\tableofcontents
\clearpage
\listoffigures
\clearpage
\listoftables

\begin{acknowledgements}
\end{acknowledgements}

\end{frontmatter}

\pagestyle{headings}

% (Optional) \part{First Part}

\chapter{Introduction}

For steam turbines, there are stators and rotors in there. These 
stators effectively act as nozzles to accelerate high pressure steam 
to supersonic speeds for use in the rotor sections of the turbine. 

For impulse turbines, the rotors do not experience pressure drop across
them (in an ideal situation), but the steam velocity definitely decreases.

The more difficult part is calculating the pressure, enthalpy and entropy 
of the steam flow out of the nozzle, especially after choked flow. What 
is especially challenging computationally is computing where the shocks 
will form. 

Where shocks form, over a constant area, we
consider the mass, energy and momentum balances:

\begin{equation*}
	\dot{m} = \rho_1 v_1 A_{shock}= \rho_2 v_2 A_{shock}
\end{equation*}

The mass balance can be simplified to become:

\begin{equation*}
	 \rho_1 v_1 = \rho_2 v_2 
\end{equation*}

\begin{equation*}
	h_0 = \frac{v^2_1}{2} + h_1 = \frac{v^2_2}{2} + h_2
\end{equation*}

\begin{equation*}
	A(P_1 -P_2) = \dot{m} (v_2 - v_1)
\end{equation*}

Fanno lines fulfil the mass and energy balance 
\citep{cengel2002thermodynamics} whereas Rayleigh lines fulfil the 
mass and momentum balances \citep{cengel2002thermodynamics}. This is on page 
861 of Cengel's thermodynamics, chapter 17 of the 8th edition.

For normal shocks, we generally expect:

\begin{equation*}
	s_{out} > s_0
\end{equation*}

However, this inequality does not help us solve the equation necessarily.
We need to plot the Fanno Lines and Rayliegh lines out in order to find out 
where the areas lie.


\chapter{Problem statement and solution sequence}

\section{How this fits in to the larger C-D nozzle picture}

Now, the problem statement for a nozzle (stator):
\begin{itemize}
	\item We have two control volumes between nozzle/rotor pair
	\item CV-1 before nozzle has known (p,h) state, and stagnation state 
	\item CV-2 after the rotor, has known (p,h) state
	\item steam goes from CV-1 to throat, accelerates 
	\item steam accelerates further or goes through shock before diverging nozzle exit 
	\item steam goes under no pressure decrease in impulse turbine, only velocity decrease 
	\item steam isentropically returns to stagnation state before entering CV-2
	\item mass flowrate is determined iteratively after thermodynamics
\end{itemize}

The shock happens specifically in the diverging part of the nozzle where 
flow is accelerated to supersonic speeds and then undergoes a shock where 
the flow is suddenly subsonic. For the case of choked flow, mass 
flowrate is known and fixed because the thermodynamic properties at 
the inlet are known.

So in this case, our mass flowrate should be known prior, because of 
the presence of choked flow, which should be for sonic conditions. Moreover,
the stagnation enthalpy and pressure before the shock should also be known.

\section{Knowns}
\begin{itemize}
	\item $\dot{m}$, choked mass flowrate
	\item stagnation enthalpy and entropy $h_0, s_0$
	\item $s_0 = s_1$ as we expect isentropic flow to that before shock
\end{itemize}
\section{Steps}

We still need to solve the mass, energy and momentum balances iteratively.
The standard way is to draw the Fanno and Rayleigh lines on a h-s diagram 
and look for its intersection. I'm not quite sure how else to do it.


\section{Rayleigh Lines}

For Rayleigh lines, we must satisfy
\begin{equation*}
	G_{shock} = \rho_1 v_1 = \rho_2 v_2 
\end{equation*}

\begin{equation*}
	A_{shock}(P_1 -P_2) = \dot{m} (v_2 - v_1)
\end{equation*}

The mass flowrate is the choked mass flowrate:

\begin{equation*}
	\dot{m} = \rho_{out} a_{throat} v_{out}
\end{equation*}

Now, we do not know what the area of the shock is, but we can eliminate 
the $A_{shock}$ from these equations:


\begin{equation*}
	(P_1 -P_2) = G_{shock} (v_2 - v_1)
\end{equation*}

Where $G$ is the mass velocity, or mass flowrate per unit cross sectional 
area.

\begin{equation*}
	G_{shock} A_{shock} = \dot{m}_{choked}
\end{equation*}

So what we eventually solve for is:
\begin{equation*}
	G_{shock} = \rho_1 v_1 = \rho_2 v_2 
\end{equation*}
\begin{equation*}
	(P_1 -P_2) = G_{shock} (v_2 - v_1)
\end{equation*}
\begin{equation*}
	(P_1 -P_2) = \rho_1 v_1 (v_2 - v_1) = \rho_2 v_2 (v_2 - v_1)
\end{equation*}

But since we have the constant mass flux or mass velocity:
\begin{equation*}
	P_1 + \rho_1 v_1^2 = P_2  + \rho_2 v_2^2
\end{equation*}

However, for Rayleigh Lines, we don't even know what the area is so as 
to determine the mass velocity.

Now, suppose we have had isentropy and a constant stagnation enthalpy,
which may help us plot things out on a (h,s) diagram. We specify a shock 
area. 
From a function perspective, we take in a mass flux, $G_{shock}$, plus 
stagnation thermodynamic states. Or rather, we specify mass flowrate, 
shock area, and then stagnation properties. From thereon, we can plot a 
Rayleigh line if we want isentropy.

Or perhaps, a better way is this. Assume $A_{shock}$ is known, this will
be user controlled. When $A_{shock}$ and $\dot{m}$ is known as inputs,
and the stagnation state $h_0, s_0$ is known, the entire input state 
is known. 

With this, we can vary $v_2$ across all values in the range 
$v_1 < v_2 < 0$ as we expect a slowdown in the velocity in the shock.

Hence, to draw the Rayleigh line:

\begin{itemize}
	\item $\dot{m}$, choked mass flowrate
	\item stagnation enthalpy and entropy $h_0, s_0$
	\item $s_0 = s_1$ as we expect isentropic flow to that before shock
	\item we specify $A_{shock}$ before starting 
	\item we calculate $G_{shock}$ 
	\item $h_0 = h_1 + v_1^2/2$
	\item $\dot{m} = \rho_1 v_1 A_{shock}$
	\item use (h,s) flash, or (p,s) flash iterate until mass balance fulfilled.
\end{itemize}

The iteration algo is

\begin{itemize}
	\item we will guess using a (h,s) flash, 
	\item expect bounds between $h_{choked}$ and $h_{saturated,water}$
	\item from guessed (h,s), $\rho_1$ and $v_1$
	\item obtain $\dot{m}$, guess again.
\end{itemize}

Note: guessing using $v$ is quite problematic, cos we need to get bounds 
of $v$

This requires a function, and we should be able to iterate our way to 
an inlet state.

Now, given this inlet state, we should know the mass flux $G_{shock}$:
\begin{equation*}
	G_{shock} = \rho_1 v_1 = \rho_2 v_2 
\end{equation*}

And we will know the LHS of this momentum balance equation:
\begin{equation*}
	P_1 + \rho_1 v_1^2 = P_2  + \rho_2 v_2^2
\end{equation*}

We can then guess sweep through values of $v_2$. In doing so, we can 
do the following steps after knowing the inlet state:

\begin{itemize}
	\item Guess $v_2$
	\item obtain $\rho_2$ from $\rho_2 = \frac{G_{shock}}{v_2}$
	\item obtain $P_2$ from: $P_2 = P_1 + \rho_1 v_1^2 - \rho_2 v_2^2$
	\item From $P_2,\rho_2$, the thermodynamic state is known, but 
		not readily obtained from steam tables
	\item Using (p,s) algorithm, we guess until $\rho_2$ is obtained, 
		from this we get $s_2$
\end{itemize}

In doing so, we have two rounds of iteration just to plot one point 
on the Fanno line, kind of a waste in my opinion. So guessing $v_2$ isn't 
quite going to cut it. 

Now, we know that the outlet pressure is known. We can use this to bound 
our equation, from low pressure all the way to outlet pressure. Doing so,
we get a pressure sweep.


\begin{itemize}
	\item Guess $h_2$ and $P_2$
	\item problem is, not all points will satisfy both momentum and 
		mass balance, there is still a need to iterate
\end{itemize}

We can try this:

\begin{itemize}
	\item $P_1 + \rho_1 v_1^2 = P_2 + G v_2$
\end{itemize}

Knowing that 
\begin{equation*}
	v_2 = G/\rho_2
\end{equation*}

\begin{itemize}
	\item $P_1 + \rho_1 v_1^2 = P_2 + \frac{G^2}{\rho_2}$
\end{itemize}

Or
\begin{equation*}
	 P_2 + \frac{G^2}{\rho_2} =\text{constant}
\end{equation*}

This means that for every density $\rho_2$, only one pressure is suitable 
for this equation. This would be a new thermodynamic state. However, it 
requires iteration to find out what this thermodynamic state is.

Aye, I'm trying to avoid iteration, but looks very hard to avoid. Seems 
that the $(p,s)$ flash is the way to go. These algorithms are unavoidable 
at least for this set of steam tables.

%But at the shockwave, energy and mass balance are:
%
%\begin{equation*}
%	h_0 = 0.5 * v_{aftershock}^2 + h_{aftershock}
%	= 0.5 * v_{beforeshock}^2 + h_{beforeshock}
%\end{equation*}
%
%\begin{equation*}
%	\rho_{beforeshock} v_{beforeshock} A_{shock} = \rho_{aftershock} v_{aftershock} A_{shock} = \dot{m}
%\end{equation*}
%
%
%\begin{equation*}
%	\dot{m} v_{beforeshock} + P_{beforeshock} A_{shock} = 
%	\dot{m} v_{aftershock} + P_{aftershock} A_{shock}
%\end{equation*}
%
%It seems as if iterative solution procedure is required to solve for 
%$A_{shock}$ and the states before and after the normal shock.
%
%The solution procedure should be as such. From choke point to before 
%the shock, the flow is isentropic. Presuming one knows $A_{shock}$, 
%a (p,s) flash is quite straightforward to determine all properties 
%of the steam before the shockwave.
%
%So if guessing $A_{shock}$, one can look at each of the balances 
%and see if they tally.
%
%\begin{itemize}
%	\item Guess $A_{shock}$
%	\item All properties before shock are known
%	\item obtain expression for mass velocity $(\rho v)_{aftershock}$
%	\item it seems iteration will also need to be done after this if we guess $A_{shock}$
%\end{itemize}
%
%Now, we may not know the $A_{shock}$ itself, but we can eliminate this 
%variable from our equations,
%
%Given $\dot{m} = \rho v A$
%
%\begin{equation*}
%	\rho v^2_{beforeshock} A_{shock}  + P_{beforeshock} A_{shock} = 
%	\rho v^2_{aftershock} A_{shock} + P_{aftershock} A_{shock}
%\end{equation*}
%
%We can eliminate $A_{shock}$ to obtain the momentum balance and mass 
%balance combined:
%
%\begin{equation*}
%	(\rho v^2)_{beforeshock}   + P_{beforeshock}  = 
%	(\rho v^2)_{aftershock}  + P_{aftershock} 
%\end{equation*}
%\begin{equation*}
%	h_0 = 0.5 * v_{aftershock}^2 + h_{aftershock}
%	= 0.5 * v_{beforeshock}^2 + h_{beforeshock}
%\end{equation*}
%
%Now, supposing that we combine these equations,
%\begin{equation*}
%	(\rho v^2)_{beforeshock} + 2 (\rho h)_{beforeshock} =
%	\rho_{beforeshock} v_{aftershock}^2 + 2 \rho_{beforeshock} h_{aftershock}
%\end{equation*}
%
%We then subtract this from both sides, to eliminate $\rho v^2$ before the 
%shock
%\begin{align*}
%	P_{beforeshock} - 2 (\rho h)_{beforeshock} =
%	P_{aftershock} - 2 \rho_{beforeshock} h_{aftershock} + \\
%	v^2_{aftershock} (\rho_{beforeshock} - \rho_{aftershock})
%\end{align*}
%
%Yeah this doesn't really work.
%
%Now, this equation is challenging to solve because there are two variables 
%to guess, $(p,h)$. To solve this, we guess $v_{aftershock}$:
%
%\begin{itemize}
%	\item guess $v_{aftershock}$
%	\item determine $h_{aftershock}$ using $h_0 = 0.5 * v_{aftershock}^2 + h_{aftershock}$
%	\item thereafter we have $ (\rho v^2)_{beforeshock}   + P_{beforeshock}  = (\rho v^2)_{aftershock}  + P_{aftershock} $
%	\item in which we guess $P$ and take $\rho(p,h)$ to be the density after shock
%	\item there may or may not be a solution for this depending on $v_{aftershock}$
%\end{itemize}
%
%Now, it's a little hard to guess where this will go, given that we 
%don't even know what $v_{beforeshock}$ will be if we don't have the area.
%
%This Fanno line problem is quite problematic!
%
%So we fix stagnation enthalpy and mass flowrate. We then vary the velocity 
%to find $h_1$, this will be our first thermodynamic variable. 
%However, there is one more degree of freedom. 
%
%From the stagnation enthalpy down to the Fanno line, is it constant 
%enthalpy?
%
%Well, supposing we specify our mass flow, $\rho$ will be determined with 
%fixed velocity, which then fully determines our thermodynamic state.
%
%Computationally, this would require iteration backwards for the steam 
%tables. It is perhaps better to find a (p,h) type equation, where pressure 
%is predetermined at the Fanno line.
%
%Honestly, I don't see an easy way to solve this.
%
%Based on the (h-s) graphs, is going through the different entropy 
%values. As seen from the balance, there will likely be two different values 
%of enthalpy that satisfy the Fanno line equations, for a given entropy 
%value.
%
%I possibly want to make an app that plots these fanno lines in real time.






\printbibliography

% \appendix
\end{document}
